\documentclass[11pt]{article}
\usepackage[a4paper,margin=26mm]{geometry}
\usepackage{amsmath,amssymb,amsthm,booktabs,array}
\usepackage[T1]{fontenc}
\usepackage[colorlinks=true,linkcolor=blue,citecolor=blue,urlcolor=blue]{hyperref}
\newtheorem{theorem}{Theorem}
\newtheorem{corollary}[theorem]{Corollary}
\title{Equal products of $m!/m^m$ over distinct parameter sets}
\author{Haobo Ma\thanks{ChronoAI Pte Ltd, Singapore.}\and
Wenlin Zhang\thanks{National University of Singapore; ChronoAI Pte Ltd, Singapore.}}
\date{November 12, 2026}
\begin{document}
\maketitle
\begin{abstract}
We answer Problem 5.6 of Aivazidis, Loukaki and M\"uller in the negative
(\emph{J. Group Theory} \textbf{27} (2024), p.~1281): the product of the weights $f(m)=m!/m^m$
does not determine a finite set of distinct integers greater than one.
The same example answers arXiv:2209.00986v1, Problem 12.6, and v2, Problem 5.6.
The two disjoint eighteen-element sets have equal products and equal parameter
sums. Their equality follows from the adjacent-weight identity
$f(n+1)/f(n)=(n/(n+1))^n$ and an explicit prime-exponent table.
For every cardinality at least nineteen, infinitely many rational values
have two distinct representations with equal cardinalities and equal sums.
\end{abstract}

\section{The question and the counterexample}
For an integer $m\geq2$, set $f(m)=m!/m^m$. For a finite set
$S\subseteq\{2,3,\ldots\}$, write
\[
 F(S)=\prod_{m\in S}f(m),\qquad F(\varnothing)=1.
\]
Aivazidis, Loukaki and M\"uller ask, ``Does Theorem 4.6 generalise?''
In particular, they ask whether equality between two such products
forces equality of the parameters after rearrangement
\cite[Problem~5.6, p.~1281]{ALM}. Equivalently, is the map $S\mapsto F(S)$
injective? The same question appears in their preprint
\cite[arXiv:2209.00986v2, Problem~5.6, p.~37]{ALMpre}. Its first version, dated September 2,
2022, poses the question without the distinctness restriction as
Problem~12.6 on p.~38 \cite[arXiv:2209.00986v1]{ALMv1}; the example below has distinct
parameters even across its two sides.

\begin{theorem}\label{thm:counterexample}
The disjoint sets
\begin{align*}
 A={}&\{10,17,21,23,24,27,38,45,46,48,56,65,68,77,80,84,95,141\},\\
 B={}&\{11,16,20,22,25,26,39,44,47,49,57,64,69,76,81,85,94,140\}
\end{align*}
satisfy $F(A)=F(B)$. In particular, the answer to Problem~5.6 is negative
even when both sides have eighteen parameters and all thirty-six parameters
are distinct. Moreover, $\sum_{a\in A}a=\sum_{b\in B}b=965$.
\end{theorem}

\begin{proof}
Define the two sets of starting indices
\begin{align*}
 U&=\{16,20,22,26,44,64,76,94,140\},\\
 V&=\{10,24,38,46,48,56,68,80,84\}.
\end{align*}
For an integer $n\geq2$, put
\begin{equation}\label{eq:adjacent}
 r(n)=\frac{f(n+1)}{f(n)}
 =\frac{(n+1)!}{(n+1)^{n+1}}\frac{n^n}{n!}
 =\left(\frac{n}{n+1}\right)^n.
\end{equation}
This identity already occurs in the proof of the source's prime-side
uniqueness theorem \cite[Theorem~4.6, p.~1263]{ALM}; the multiplicative
relation below uses distinct adjacent pairs.
For a prime $p$ and a positive rational number $x$, let $v_p(x)$ be the
exponent of $p$ in its prime factorization; this exponent can be negative.
For $S=U,V$, define
\[
 E_p(S)=\sum_{n\in S} n\bigl(v_p(n)-v_p(n+1)\bigr)
       =v_p\!\left(\prod_{n\in S}r(n)\right).
\]
Factoring the displayed integers gives the following table:
\begin{center}
\small
\setlength{\tabcolsep}{4.3pt}
\begin{tabular}{@{}lrrrrrrrrrr@{}}
\toprule
$p$ &2&3&5&7&11&13&17&19&23&47\\
\midrule
$E_p(U)$ &1150&$-326$&$-42$&44&$-10$&$-38$&$-16$&$-18$&$-22$&$-46$\\
$E_p(V)$ &1150&$-326$&$-42$&44&$-10$&$-38$&$-16$&$-18$&$-22$&$-46$\\
\bottomrule
\end{tabular}
\end{center}
No other prime divides any of $n$ or $n+1$ for $n\in U\cup V$.
For example, $140=2^2\cdot5\cdot7$ and $141=3\cdot47$, so the
term $n=140$ contributes $280, -140, 140, 140, -140$ at primes
$2,3,5,7,47$, respectively. Thus the two products of adjacent ratios are
both
\begin{equation}\label{eq:certificate}
 \prod_{n\in U}r(n)=\prod_{n\in V}r(n)
 =\frac{2^{1150}7^{44}}
 {3^{326}5^{42}11^{10}13^{38}17^{16}19^{18}23^{22}47^{46}}.
\end{equation}
For any set $S$ of integers, write $S+1=\{n+1:n\in S\}$.
All elements of $U\cup V$ are distinct and even. Hence their adjacent
pairs $\{n,n+1\}$ are pairwise disjoint. In particular,
\[
 A=V\cup(U+1),\qquad B=U\cup(V+1).
\]
Substituting \eqref{eq:adjacent} into \eqref{eq:certificate} and multiplying
by $F(U)F(V)$ proves $F(A)=F(B)$. The disjointness and cardinalities follow
from the eighteen disjoint adjacent pairs. Every parameter exceeds one,
and $10\in A\setminus B$, so the sets are different.
Finally, each side chooses nine upper and nine lower endpoints of the same
eighteen pairs. Both sums are therefore
$\sum_{n\in U\cup V}n+9=956+9=965$.
\end{proof}

\begin{corollary}\label{cor:extensions}
For every integer $k\geq18$, there are distinct $k$-element sets
$A_k,B_k\subseteq\{2,3,\ldots\}$ with $F(A_k)=F(B_k)$ and
$|A_k\mathbin{\triangle}B_k|=36$, where $\triangle$ denotes symmetric
difference. Their parameter sums are equal. More generally, any finite set
$C$ disjoint from $A\cup B$ can be adjoined to both sides, giving equal
products and equal sums $965+\sum_{c\in C}c$.
\end{corollary}
\begin{proof}
Choose any $(k-18)$-element set $C$ disjoint from $A\cup B$, and take
$A_k=A\cup C$ and $B_k=B\cup C$. Both products are multiplied by $F(C)$,
and their symmetric difference is $A\cup B$. Their sums both increase by
$\sum_{c\in C}c$. For $k>18$ these two sets share $C$; disjointness is
asserted for the original pair, not for its common-parameter extensions.
\end{proof}

\begin{corollary}\label{cor:infinite}
For each integer $k\geq19$, infinitely many positive rational numbers are
values of $F$ at two different $k$-element sets with equal parameter sums.
\end{corollary}
\begin{proof}
Fix a $(k-19)$-element set $D$ disjoint from $A\cup B$. For each integer
$c>\max(A\cup B\cup D)$, adjoin $D\cup\{c\}$ to both $A$ and $B$.
The common product is $F(A)F(D)f(c)$. Equation~\eqref{eq:adjacent} gives
$0<f(c+1)<f(c)$, so these common products are pairwise different.
\end{proof}

\begin{corollary}\label{cor:comparable}
Suppose $a_1<\cdots<a_k$ and $b_1<\cdots<b_k$ are integers greater than
one, with $a_i\leq b_i$ for every $i$. Then
$\prod_i f(a_i)\geq\prod_i f(b_i)$, with equality exactly when
$a_i=b_i$ for every $i$.
\end{corollary}
\begin{proof}
Equation~\eqref{eq:adjacent} makes $f$ positive and strictly decreasing
on the integers greater than one. Thus every factor on the left is at
least its corresponding factor on the right, and any strict parameter
inequality gives a strict product inequality.
\end{proof}

\section{Source coverage and mathematical evidence}
The prime-restricted uniqueness theorem
\cite[Theorem~4.6, p.~1263]{ALM} remains distinct from Problem~5.6:
it requires one side to consist of distinct primes and permits arbitrary
integer parameters on the other side. Both sets in
Theorem~\ref{thm:counterexample} contain composite integers, so they do not
meet that prime-side hypothesis. The example concerns the arithmetic
question and imposes no group-realization condition. It also shows that
prescribing both cardinality and parameter sum does not restore unrestricted
uniqueness; this strengthened uniqueness assertion was not posed in the source.

The source's majorisation comparisons \cite[Appendix~A, pp.~1281--1283]{ALM}
have a different hypothesis. In decreasing order, the first partial sums
are $141>140$, but the first nine give $714<715$. Thus the two vectors
are incomparable in weak majorisation. The complete partial-sum lists are
supplied by \texttt{verify.py}.

\begin{center}
\small
\begin{tabular}{@{}>{\raggedright\arraybackslash}p{0.31\linewidth}>{\raggedright\arraybackslash}p{0.62\linewidth}@{}}
\toprule
Source item & Coverage\\
\midrule
Problem~5.6, p.~1281 & Answered: no. Theorem~\ref{thm:counterexample} gives disjoint eighteen-element sets with equal products and equal sums.\\
Problem~12.6 in v1, p.~38 & Answered: no; the same example also permits the original unrestricted lists.\\
Theorem~4.6, pp.~1263--1265 & Prior proof cited: uniqueness with one side consisting of distinct primes.\\
Appendix~A, pp.~1281--1283 & Prior comparison results cited; the witness is incomparable in weak majorisation.\\
Problem~5.1, p.~1280 & Not covered: realization by finite groups and subgroups.\\
Problem~5.2, p.~1280 & Not covered: normal-extension inequality and its equality cases.\\
Problem~5.3, p.~1280 & Not covered: equality for direct products.\\
Problem~5.4, p.~1280 & Not covered: infinite-group formulation.\\
Problem~5.5, p.~1280 & Not covered: comparison with commuting probability and equality cases.\\
Problem~5.7, p.~1281 & Not covered: comparison for nested subgroups.\\
Problem~5.8, p.~1281 & Not covered: cyclic prime-power subgroup attaining the minimum.\\
Uniqueness at fixed size and sum & Not asserted in the source; fails for the displayed eighteen-element sets of sum $965$.\\
Extensions not posed there & Proved: Corollaries~\ref{cor:extensions}--\ref{cor:comparable}.\\
\bottomrule
\end{tabular}
\end{center}

\subsection*{Verification scope}
The pinned Lean declaration
\href{https://github.com/the-omega-institute/trureturing/blob/736015507f370926a4b39c6a5ca4f090fc50cae3/D5/S3/FiniteGroups/FactorialProductNonunique.lean}{\texttt{FactorialProductNonunique.result}}
proves $\neg\texttt{claim}$, where \texttt{claim} asserts that for all finite
subsets $A,B$ of the natural numbers whose elements exceed one, equality of
the factorial-weight products implies $A=B$. Its witness is exactly the
two sets above. Its weight is
\href{https://github.com/the-omega-institute/trureturing/blob/736015507f370926a4b39c6a5ca4f090fc50cae3/D5/S3/FiniteGroups/TransversalProbabilityMonotone.lean}{\texttt{TransversalProbabilityMonotone.orbitWeight}},
defined by $(m!:\mathbb Q)/(m:\mathbb Q)^m$. This is $f(m)$ on the
stated domain. This declaration does not state the equal-sum or
arbitrary-cardinality conclusions. The adjacent-ratio argument and
Corollaries~\ref{cor:extensions}--\ref{cor:comparable} are written proofs. The accompanying
\texttt{verify.py} checks the explicit factorial products and the exponent
table, parameter sums and majorisation partial sums using exact integer and
rational arithmetic. It also tests finitely many common-parameter extensions;
the conclusions for all cardinalities and infinitely many values follow
from the written proofs. The weight's dependency file also contains a
normal-extension probability theorem, which is not needed for this argument.

\begin{thebibliography}{3}
\bibitem{ALM}
S.~Aivazidis, M.~Loukaki and T.~W.~M\"uller,
\emph{On the common transversal probability},
J.~Group Theory \textbf{27} (2024), no.~6, 1233--1285.
\href{https://doi.org/10.1515/jgth-2024-0030}{doi:10.1515/jgth-2024-0030}.

\bibitem{ALMpre}
S.~Aivazidis, M.~Loukaki and T.~W.~M\"uller,
\emph{On the common transversal probability in finite groups},
\href{https://arxiv.org/abs/2209.00986v2}{arXiv:2209.00986v2},
May 31, 2023.

\bibitem{ALMv1}
S.~Aivazidis, M.~Loukaki and T.~W.~M\"uller,
\emph{On the common transversal probability in finite groups},
\href{https://arxiv.org/abs/2209.00986v1}{arXiv:2209.00986v1},
September 2, 2022; Problem~12.6, p.~38.
\end{thebibliography}

\bigskip
\noindent\footnotesize
\copyright\ 2026 Haobo Ma and Wenlin Zhang.
Licensed under \href{https://creativecommons.org/licenses/by/4.0/}{Creative
Commons Attribution 4.0 International (CC BY 4.0)}.
\end{document}
